\documentclass[11pt]{article}
\usepackage[utf8]{inputenc}
\usepackage[T2A]{fontenc}
\usepackage[russian]{babel}
\usepackage{amsmath,amssymb}
\usepackage[margin=2.5cm]{geometry}
\usepackage{hyperref}

\title{Разделение секрета (Shamir) и MPC между роботами}
\author{}
\date{11 сентября 2026}

\begin{document}
\maketitle

\section{Примитив: Shamir Secret Sharing}

Схема Шамира \cite{shamir1979} разбивает секрет $s$ на $n$ долей так, что любые $t$
долей (порог) полностью восстанавливают $s$, а любые $t-1$ долей не дают о нём
\emph{никакой} информации в информационно-теоретическом смысле. Строится случайный
многочлен степени $t-1$:
\[
  f(x) = s + a_1 x + a_2 x^2 + \dots + a_{t-1} x^{t-1}, \qquad f(0) = s,
\]
где $a_1,\dots,a_{t-1}$ выбираются случайно. Участнику $i$ выдаётся точка
$(i,\,f(i))$. По интерполяции Лагранжа любые $t$ точек однозначно восстанавливают
$f$, а значит и $s = f(0)$; при $< t$ точках многочленов, проходящих через них,
бесконечно много — секрет не определён.

\section{Встраивание в secure aggregation}

Канонический протокол \textbf{secure aggregation} для федеративного обучения
\cite{bonawitz2017} использует Shamir SS не для самих данных, а для \emph{масок}:
\begin{enumerate}
  \item Каждый клиент $k$ маскирует свой вклад $x_k$ попарными псевдослучайными
        масками, сгенерированными через общий ключ с каждым другим клиентом (PRF).
  \item При суммировании на сервере все маски взаимно уничтожаются:
        $\sum_k (x_k + \text{маски}_k) = \sum_k x_k$ — сервер видит только сумму,
        не отдельные вклады.
  \item Если клиент отваливается посреди протокола (для флота роботов —
        частая ситуация: Wi-Fi, разряд батареи), его маска \guillemotleft
        зависает\guillemotright{} и портит сумму. Доли ключей/масок каждого клиента
        заранее разложены схемой Shamir $(t, n)$ среди остальных клиентов: при потере
        клиента оставшиеся (если их $\ge t$) совместно реконструируют недостающую
        маску и вычитают её вклад.
\end{enumerate}

\textbf{Важное следствие для проекта:} связка \guillemotleft MPC + Shamir\guillemotright{}
не новый метод — это прямое применение протокола Bonawitz et al. 2017 года, на
котором построена Google Secure Aggregation. Claim новизны на этом уровне
несостоятелен; новизна проекта должна опираться на что-то другое (атака
реконструкции, защита с сохранением полезности — см. отдельные материалы проекта).

\section{Расширения протокола}

Post-quantum вариант \cite{pqsf2024} заменяет DH-ключи маскирования на решёточную
криптографию для долгосрочной устойчивости. Верифицируемая secure aggregation
\cite{groupverifiable2025} добавляет возможность группе клиентов доказать корректность
агрегации без раскрытия вкладов — актуально при переходе от модели
\guillemotleft honest-but-curious\guillemotright{} сервера к недоверенному серверу.
Гибрид Shamir SS с 2-из-2 аддитивным secret sharing через PRF убирает прямую
попарную коммуникацию между клиентами \cite{sokseragg} — но для флота роботов это не
даёт выигрыша: физически не все пары роботов могут связаться напрямую (разные дома),
поэтому коммуникация в любом случае идёт через сервер как relay.

\section{Незакрытая угроза: byzantine-роботы}

Secure aggregation защищает от подглядывания \emph{сервером}, но не от
\textbf{испорченных клиентов}, которые честно маскируют, но пересылают вредоносные
обновления \cite{sable2023}. В исходной постановке проекта эта угроза не рассмотрена
вовсе: испорченный робот в потенциале способен исказить восстановленный глобальный
квантиль или агрегированные веса даже при полностью корректной работе Shamir-слоя.
Это самостоятельная и открытая ось для расширения C2 (см. основной security-роадмап
проекта).

\section{Тезаурус}

\begin{description}
  \item[Secret sharing (разделение секрета)] Разбиение секретного значения на доли
    между участниками так, что восстановить секрет можно только собрав определённое
    минимальное число долей.
  \item[$(t, n)$-пороговая схема] Схема разделения секрета на $n$ долей, где любые
    $t$ из них восстанавливают секрет, а меньше $t$ — не дают о нём никакой
    информации.
  \item[Интерполяция Лагранжа] Метод однозначного восстановления многочлена степени
    $t-1$ по $t$ известным точкам — математическая основа схемы Шамира.
  \item[MPC (Multi-Party Computation)] Класс протоколов, позволяющих нескольким
    недоверяющим друг другу сторонам совместно вычислить функцию от своих секретных
    входов, не раскрывая сами входы.
  \item[Secure aggregation] Протокол MPC, в котором сервер получает только сумму
    (или другую агрегированную величину) вкладов клиентов, не видя вклад каждого по
    отдельности.
  \item[PRF (Pseudo-Random Function)] Функция, генерирующая псевдослучайный выход по
    ключу — используется парами клиентов для согласованной генерации взаимоуничтожающихся
    масок.
  \item[Маскирование] Добавление к вкладу клиента случайного числа (маски) перед
    отправкой, так что при суммировании всех вкладов маски взаимно уничтожаются.
  \item[Honest-but-curious] Модель угроз, в которой участник (например, сервер)
    строго следует протоколу, но пытается извлечь дополнительную информацию из того,
    что видит.
  \item[Byzantine-угроза] Угроза от участника, который \emph{не} следует протоколу
    честно — например, отправляет заведомо испорченные (отравленные) обновления.
  \item[Верифицируемость] Возможность для участников протокола доказать корректность
    выполненного вычисления, не раскрывая исходные данные.
\end{description}

\begin{thebibliography}{9}
\bibitem{shamir1979} A.~Shamir. \emph{How to Share a Secret}. Communications of the
  ACM, 22(11), 1979.
\bibitem{bonawitz2017} K.~Bonawitz et al. \emph{Practical Secure Aggregation for
  Privacy-Preserving Machine Learning}. ACM CCS, 2017.
\bibitem{pqsf2024} \emph{PQSF: Post-Quantum Secure Privacy-Preserving Federated
  Learning}. Scientific Reports, 2024.
\bibitem{groupverifiable2025} \emph{Group Verifiable Secure Aggregate Federated
  Learning Based on Secret Sharing}. Scientific Reports, 2025.
\bibitem{sokseragg} \emph{SoK: Secure Aggregation based on Cryptographic Schemes for
  Federated Learning}. EURECOM.
\bibitem{sable2023} \emph{SABLE: Secure And Byzantine robust LEarning}.
  arXiv:2309.05395, 2023.
\end{thebibliography}

\end{document}
